\documentclass[11pt]{article}
\usepackage[margin=1in]{geometry}
\usepackage{geometry}
\setlength{\parskip}{0.6em}
\setlength{\parindent}{0pt}
\begin{document}

\hfill November 14, 2026

\vspace{0.5em}

The Editors\\
Structural Change and Economic Dynamics

\vspace{0.8em}

Dear Editors,

I am pleased to submit my manuscript entitled ``Anatomy of a Smooth Transition:
A Task-Based Macroeconomic Model of AGI Displacement with Fiscal and Social
Safety Nets'' for consideration for publication in \emph{Structural Change and
Economic Dynamics}.

The paper studies the labor transition under artificial general intelligence
(AGI) in a task-based macroeconomic model with an endogenously collapsing
reinstatement engine, nominal rigidity and the zero lower bound, a fiscal
closure, and a social-instability channel that is independent by construction.
Under full substitution, the direct automation flow persistently exceeds the
labor capture of new tasks, so within the model the market employment share
falls monotonically to zero and a high market-employment endpoint is
inadmissible; that branch is reported as a counterfactual rather than a
reachable steady state. Once public
employment is set to absorb the entire unemployed population, the global
outcome is a two-basin bifurcation between an economic crash and a smooth
transition to full public-employment absorption, and a bisection locates the
critical frontier growth that separates them.

The work fits the journal's scope directly. It is a study of technology-driven
structural change analyzed through nonlinear dynamics, including a
piecewise-smooth boundary equilibrium bifurcation and geometric singular
perturbation theory, rather than a single comparative-static equilibrium. The
task boundary between humans and AGI takes the place of the sector boundary in
the structural-change tradition, and the fiscal and social safety nets are
closed explicitly.

This manuscript is original, is the sole work of the named author, has not been
published previously, and is not under consideration elsewhere. There is no
conflict of interest. The simulation and analysis code is provided with the
submission and every numerical result in the paper is produced by a named,
re-runnable script. A separate online appendix documents the external
empirical anchoring and the extended robustness checks.

I would be glad to provide any further information you may require.

Sincerely,

\vspace{1.2em}

DUAN Fang\\
Corresponding author\\
College of Information Science and Engineering, Huaqiao University,
Xiamen, China\\
duan.fang@hqu.edu.cn

\end{document}
